\section{Attack Suite and What Makes an Attack Strong}
\label{sec:attacks}

The threat model of \S\ref{sec:threat-model} defines an adversary; this
section instantiates it. We describe the attacks we run, the scoring functions
they apply, the metric we read them by, and --- because it governs all three
--- what we take attack strength to consist of.

\subsection{What makes an attack strong}
\label{sec:attacks-strength}

The privacy literature routinely ranks attacks by adversary access level.
Under \emph{black-box} access the attacker observes only model outputs; under
\emph{white-box} access it observes weights, gradients and internal
activations. The intuition is that the latter is more dangerous because it
sees more, and from that intuition follows a widespread defensive strategy:
reduce what the adversary can observe, hide the gradients, encrypt the
updates.

Our apparatus rests on a different premise, and our results support it.
\emph{When a sample's loss is observable, attack strength is determined not by
access level but by the quality of the scoring function.}

Loss occupies a special position. A membership attack must answer a binary
question: was this sample in the training set? Anything that helps answer it
must surface in some measurable quantity, and for a model trained to minimise
a loss function that quantity is the loss itself. A sample the model trained
on is a sample whose loss the optimiser pushed down; an unseen sample is not.
Memorisation, by construction, is visible in the loss, because minimising the
loss is what training does.

Sablayrolles et al.~\cite{sablayrolles2019whitebox} turn this into a formal
result. Assuming the adversary knows the loss distributions for in-training
and out-of-training samples, there exists a transformation of the loss that is
\emph{Bayes-optimal}: it extracts all available membership information. Not
part of it --- all of it. No further observation, whether per-layer gradients,
intermediate activations or the optimiser trajectory, can take the attacker
past that bound.

White-box access retains a role, but an indirect one. The adversary does not
truly know the two distributions; it estimates them, typically by training
shadow models. Privileged access can improve that estimate, and thus move the
attacker closer to the bound --- but through better calibration, not because
gradients carry membership signal that the loss does not.

Two consequences follow, and the second is a design recommendation.

For our measurement, we concentrate experimental effort on evaluator quality
rather than on multiplying access levels: three loss transformations
(\S\ref{sec:attacks-scorers}), per-example shadow calibration, and explicit
checks that scoring treats both classes under the same formula
(\S\ref{sec:validation}). Our data confirm the choice. LiRA is the most
privileged attack we run, observing an individual client's update before
aggregation, yet when its calibration degenerates it fails to beat a plain
threshold test on raw loss.

For defence designers, hiding gradients does not protect if the loss remains
observable. A defence that reduces adversary visibility without reducing the
separability of the two loss distributions shifts the cost of the attack
rather than its feasibility: it makes estimation more laborious, not the
conclusion less reachable. This is false assurance, and it must be
distinguished from defences that act on separability --- differential privacy
among them --- which reduce the signal rather than obscure it.

\subsection{Two observation surfaces, three attacks}
\label{sec:attacks-surfaces}

The attacks we run do not all operate on the same artefact, and the
distinction bounds what each can conclude.

Yeom et al.'s loss threshold~\cite{yeom2018privacy} and a calibrated
shadow-model attack~\cite{shokri2017membership} both observe the
\emph{aggregated global model}, available to any participant. A positive
detection concludes that a session was a training member somewhere in the
federation, never at which client.

LiRA~\cite{carlini2022membership} observes an \emph{individual client's per-round
update} as submitted for aggregation. It therefore requires aggregator
privilege, and a positive detection attributes membership to a specific
operator.

These are not three points on one scale of strength. They are two independent
surfaces testing two different claims, and we report them separately. LiRA is
our primary attack because it matches the aggregator's real privilege under
our threat model, not because it subsumes the other two.

\subsection{Three scoring functions}
\label{sec:attacks-scorers}

Every attack we run derives its judgement from the same observable quantity,
the target model's reconstruction error on the sample. They differ in the
transformation applied before deciding, and we implement the three the
literature distinguishes.

\textbf{Raw loss.} The untransformed mean squared error. The assumption is
that a sample seen in training is reconstructed better. The limitation is
known and dominant: a sample's loss depends chiefly on how intrinsically
difficult that sample is, not on whether it was trained on. An atypical record
has high loss regardless, which confounds membership with difficulty.

\textbf{Log-transformed loss.} The natural logarithm of the loss. Its purpose
is variance stabilisation. Raw loss distributions are strongly right-skewed,
with a long tail of large errors, whereas LiRA assumes shadow losses are
approximately Gaussian. The assumption is not cosmetic: LiRA estimates a mean
and a variance per sample, and on a skewed distribution those estimates
describe nothing. On our data the violation is measured and severe --- a
Jarque--Bera test on raw losses returns values six orders of magnitude above
the rejection threshold, and the log transformation reduces the departure from
normality by four to five orders.

\textbf{Bayes-optimal score.} The transformation derived by Sablayrolles et
al.~\cite{sablayrolles2019whitebox}, theoretically the best possible mapping
for this hypothesis test. It is the upper bound of
\S\ref{sec:attacks-strength} made concrete: no further manipulation of the
loss, and no access to internal gradients, improves on it.

Implementing all three is not redundancy. It is how we separate the effect of
calibration from the effect of signal: if the three agree, the limit lies in
the data; if they diverge, the limit lies in the transformation, and we know
which one.

\subsection{Why AUC is not enough}
\label{sec:attacks-metrics}

A membership attack produces, for each sample, a score ranking candidates from
most to least likely member. AUC-ROC measures that ranking as a whole: the
probability that a randomly chosen member outranks a randomly chosen
non-member. An AUC of $0.5$ is random ordering, $1.0$ perfect separation.

This question corresponds to no real adversary. Carlini et
al.~\cite{carlini2022membership} argue the point: an attacker need not classify
every sample correctly and does not care about average accuracy. What the
attacker wants is to identify \emph{some} members with confidence. An
adversary who, across ten thousand sessions, states with near certainty that
fifty were in a given operator's training set has breached fifty people's
privacy; guessing at random on the remaining 9,950 neither helps nor hurts
that claim.

AUC averages exactly this away. It integrates over every threshold, weighting
equally the high-false-positive regime --- where the attacker accuses half the
dataset to catch half of it, which is not an attack --- and the
low-false-positive regime, where the attacker accuses few samples and is
right. An attack can therefore identify one per cent of members with certainty
while showing an AUC barely distinguishable from $0.5$, because that one per
cent contributes almost nothing to the integral.

We therefore report true-positive rate at fixed low false-positive rates
($0.1\%$, $1\%$, $5\%$) and treat it as primary for LiRA. The question it asks
is the operative one: \emph{if the attacker accepts being wrong at most once in
a hundred, how many members can be identified?} It measures the tail of the
score distribution, which is where harm concentrates.

Two points we keep explicit. The chance level for this metric is
$\mathrm{TPR} = \mathrm{FPR}$, not $0.5$: an attacker guessing at random and
accusing $1\%$ of samples catches $1\%$ of members. And the two metrics can
diverge --- the region where they diverge is precisely where the attack is
dangerous --- so reporting both is not redundancy.

We nonetheless retain AUC, for two practical reasons: it is the metric against
which the empirical literature compares, and the equivalence test supporting
our null result is formulated on it.
